\section{Preliminaries}\label{sec:prelim}

\paragraph{Numerical semigroups.} Here $\N=\{0,1,2,\dots\}$, and an interval $[a,b]$ means the set of integers $n$ with $a\le n\le b$ (likewise for half-open intervals), except on the circle $\mathbb R/\Z$ in Section~\ref{sec:family}. A numerical semigroup is a submonoid $S$ of $(\N,+)$ with finite complement. Its elements not in $S$ are its gaps, its largest gap is the Frobenius number $\Fr(S)$, and $c(S)=\Fr(S)+1$ is its conductor; for $S=\N$ we put $\Fr(S)=-1$ and $c(S)=0$. For a positive element $m$ of $S$, the Ap\'ery set of $S$ with respect to $m$ is $\Ap(S;m)=\{w_r:r\in\Z/m\Z\}$, where $w_r$ is the least element of $S$ congruent to $r$ modulo $m$; an integer $n\ge0$ lies in $S$ if and only if $n\ge w_{n\bmod m}$, and $\Fr(S)=\max_r w_r-m$ \cite{RosalesGarciaSanchez2009}.

\begin{definition}[Mass]\label{def:mass}
The \emph{mass} of a numerical semigroup $S$ is
\[
\M(S)=\sum_{x\notin S}\bigl(x-\max\{s\in S: s<x\}\bigr),
\]
the sum over the gaps of $S$ of the distance down to the preceding element of $S$.
\end{definition}

A desert, that is, a maximal run of consecutive gaps \cite{NumericalSgps} (for two generators see \cite[Sect.~6]{Moree2014}), of length $L$ contributes $1+2+\dots+L=L(L+1)/2$, and so does the sum of the distances up to the next element. For a finite set $X\subset\Z$ write $f(X)$ for the number of pairs $s\le t$ with $[s,t]\subseteq X$; if $X$ is a union of maximal runs of lengths $L_i$, then $f(X)=\sum_i L_i(L_i+1)/2$. Thus $\M(S)=f(\N\setminus S)$.

\paragraph{Circular gaps.} We use one statistic of arithmetic progressions on a circle. For coprime integers $p\ge1$ and $r$, and $1\le h\le p$, let $g_1,\dots,g_h$ be the circular spacings of the points $0,r,\dots,(h-1)r$ in $\Z/p\Z$, and put
\[
\phi_{p,r}(h)=\sum_{i=1}^h\binom{g_i}2,\qquad\Phi(p,r)=\sum_{h=1}^p\phi_{p,r}(h),
\]
with $\phi_{p,r}(h)=0$ for $h>p$. Then $\phi_{p,r}(1)=\binom p2$, $\phi_{p,r}(p)=0$, and $\phi_{p,r}$ depends only on $r\bmod p$. We write $\M(p,q)=\M(\langle p,q\rangle)$ for coprime $p,q\ge1$.

\paragraph{Sumsets and undershoot.} Throughout, $A_0=\{0=a_0<a_1<\dots<a_k=d\}$ is a set of integers with $\gcd(A_0)=1$; dividing by the gcd reduces the general case to this one. For $h\ge0$, $hA_0$ is the set of sums of $h$ elements of $A_0$, that is, of at most $h$ elements of $A=A_0\setminus\{0\}$. For an integer $T\ge0$ the \emph{undershoot} is
\[
D_h(T)=T-\max\{n\in hA_0: n\le T\},
\]
and $\Tot(h)=\sum_{T=0}^{hd}D_h(T)$ is the total undershoot. We write $S_1=\langle a_1,\dots,a_k\rangle$ and $S_2=\langle d-a_0,\dots,d-a_{k-1}\rangle$ for the two numerical semigroups of $A_0$, with conductors $c_1$ and $c_2$.

\paragraph{The homogenized semigroup.} Let $\mathcal S\subset\N^2$ be the semigroup generated by the points $(a,d-a)$, $a\in A_0$. Then $n\in hA_0$ if and only if $(n,hd-n)\in\mathcal S$; the semigroup ring $\Bbbk[\mathcal S]$ over a field $\Bbbk$ is the homogeneous coordinate ring of the projective monomial curve $C_A$. Let $G$ be the group generated by $\mathcal S$.

\begin{theorem}[Cavaliere and Niesi {\cite[Lemma~4.3, Thm.~4.6]{CavaliereNiesi1983}}]\label{thm:cm}
The ring $\Bbbk[\mathcal S]$ is Cohen--Macaulay if and only if $\mathcal S=G\cap(S_1\times S_2)$. Equivalently, for every $h\ge1$,
\[
hA_0=I_h:=\{n\in[0,hd]: n\in S_1,\ hd-n\in S_2\}.
\]
\end{theorem}

We call $A_0$ \emph{CM} when these conditions hold, and $I_h$ the \emph{intersection form} at level $h$; in general $hA_0\subseteq I_h$. The same statement is \cite[Prop.~2.6]{GimenezGonzalezSanchez2023}; for general simplicial affine semigroups see \cite[Thm.~(2.6)]{GotoSuzukiWatanabe1976} and \cite[Cor.~4.4]{TrungHoa1986}.

\paragraph{Sumsets regularity.} Following \cite[Def.~1.3]{GimenezGonzalezSanchez2023}, the sumsets regularity $\sigma(A)$, a Hilbert-function index of $C_A$ and not its Castelnuovo--Mumford regularity, is the least $h_0\ge0$ such that for every $h\ge h_0$ the interval $[c_1,hd-c_2]$ is nonempty and
\[
hA_0=\bigl(S_1\cap[0,c_1-2]\bigr)\ \sqcup\ [c_1,hd-c_2]\ \sqcup\ \bigl(hd-(S_2\cap[0,c_2-2])\bigr).
\]
The requirement that the central interval be nonempty, that is $hd\ge c_1+c_2$, matters: without it the decomposition can hold at a smaller level (for $\{0,3,4,6\}$ it holds already at $h=1$ with an empty central interval). By \cite[Thm.~3.1, Prop.~3.2]{GimenezGonzalezSanchez2023}, $\sigma(A)$ is the maximum of the Hilbert regularity index of $\Bbbk[\mathcal S]$ and $\lceil(c_1+c_2)/d\rceil$, and equals $\lceil(c_1+c_2)/d\rceil$ when $A_0$ is CM. Gimenez and Gonz\'alez-S\'anchez assume that $A_0$ has at least four elements; for three elements $C_A$ is a plane curve of degree $d$, hence arithmetically Cohen--Macaulay with $\reg\Bbbk[\mathcal S]=d-1$, which is also the value of \eqref{eq:reg-cm} below, and the accompanying code checks the other facts cited from \cite{GimenezGonzalezSanchez2023} on the three-element sets of Table~\ref{tab:checks}.

We use four elementary facts throughout.

\begin{lemma}[Elementary facts]\label{lem:elementary}
\begin{enumerate}
\item[(i)] If $T\ge hd$ then $D_h(T)=T-hd$.
\item[(ii)] For parts with greatest common divisor $g$, the undershoot of a target $T$ is $(T\bmod g)+g\,D_h(\lfloor T/g\rfloor)$ for the reduced set.
\item[(iii)] Removing an element $x$ with $0<x<hd$ from $hA_0$ strictly increases $\Tot(h)$.
\item[(iv)] On each maximal run of $L$ consecutive non-members of $hA_0$ inside $[0,hd]$, the distances down to the preceding member and the distances up to the following member both add up to $L(L+1)/2$.
\end{enumerate}
\end{lemma}

\begin{proof}
(i) Every element of $hA_0$ is at most $hd$, and $hd\in hA_0$. (ii) The elements of $hA_0$ are the multiples $gn$ with $n$ in the $h$-fold sumset of the reduced set, so the largest one not above $T$ is $g$ times the largest one not above $\lfloor T/g\rfloor$. (iii) The undershoot at $x$ becomes positive and no undershoot decreases. (iv) Both sums are $1+2+\dots+L$.
\end{proof}
